\section{Week 16 CONTINUE--VERIFY--MATRIX--CRT Stop Boundary}

The continuation verdict is \status{stop-verify-matrix-crt}, not
\status{go-verify}.  The requested
\procname{verify_matrix_crt_mode2_fromcrt_freeze_exact} theorem was not
proved, and neither V-3/V-4 nor the desired output equality was introduced as
a premise.

\subsection*{Preserved baseline and actual call chain}

The pre-continuation aggregate was copied to
\breakablett{logs/verify-all-before-verify-matrix-crt.log}.  It contains 82 fresh
compile lines and the terminal result
\texttt{RESULT PASS authored-targets=82 cache=-no-eco}; its SHA-256 is
\breakablett{46e7dac8e442c820f746139a164c8bc00d6af17b7ad25cbd5d195507fddae03c}.
Fresh extraction confirms that actual \procname{_verify_matrix_crt} calls
\procname{_polyvec_ntt}, \procname{_polymat_pointwise_acc},
\procname{_polyvec_invntt}, \procname{_polyveck_poly_fromcrt}, and
\procname{_polyvec_freeze2q} in that order.  Mode 2 fixes two rows and four
columns.

The generated from-CRT body lifts words (0\ldots255) using the parity of
\texttt{high xor wprime}, lifts words (256\ldots511) using the parity of
\texttt{high}, and then applies the generated \procname{__freeze2q} to all 512
active words.  Reading those bodies establishes the proof target; it is not a
substitute for a checked Hoare theorem.

The post-continuation single-writer aggregate is preserved in
\breakablett{logs/verify-all-week16-verify-matrix-crt.log}.  It again has exactly
82 fresh compile lines and one terminal result, with SHA-256
\breakablett{4cd64e5a656be82710bca1410c4d19403a3c661d6b91b0319a0ea8f7c91646da}.

\subsection*{Exact missing leaves}

The first missing procedural leaf in program order is
\procname{verify_matrix_ntt_acc_mode2_cols4_correct}.  It must identify the
two returned inverse-NTT rows with the four-term negacyclic polynomial row
products.  Existing reusable arithmetic theorems fix the KeyGen shape to two
rows and three columns, so they cannot instantiate this Verify result.

The first concrete pure-algebra subleaf encountered beneath it is
\procname{rq_mul_coeff_foldr_to_bigi}: the \texttt{Rq.\&*} coefficient
definition must be normalized from its
\texttt{foldr (iota\_ 0 256)} form to a finite negacyclic sum.  The next absent
leaf is \procname{full_ntt_montgomery_spectral_action}, which must prove

\[
  \operatorname{mont}\!\left(\operatorname{full\_invntt}
    (\widehat a\,\operatorname{full\_ntt}(p)R^{-1})\right)
  = \operatorname{full\_invntt}(\widehat a)\mathbin{\&*}p.
\]

Its finite-sum proof also needs the missing 256-term odd-root orthogonality
identity.  The NTT artifact proves imperative-to-closed-form transform
correctness, not this spectral action.

The later procedural sibling is
\procname{verify_crt_freeze_mode2_word_exact}.  It must prove the two parity
lifts, the 512-word freeze, and the inactive-tail frame for the actual loops.
A direct invariant attempt failed fresh compilation with
\texttt{nothing to introduce}; the non-compiling block was removed.  Only
after both procedural leaves are checked can the requested combined theorem
be composed.

\subsection*{Frozen scope}

The continuation changes none of the existing V-1, V-2, V-5, V-6, norm, or
tail theorems.  It does not enter KeyGen NTT, Sign challenge semantics,
codecs, parsing, public APIs, probability, termination, or security.  The
82-target aggregate and all extraction, premise, hole, axiom, debug, scope,
LaTeX, and source-drift gates pass again, but no missing equality is treated
as an assumption.  Therefore the final decision remains
\status{stop-verify-matrix-crt}.
